\section{Architecture decision: Lambda or Kappa}

\subsection{The two options}
Both architectures keep raw events in an append-only store and build views from them \cite{kleppmann2017}. They differ in how
they compute and recompute those views.

\textbf{Lambda} has three layers \cite{marz2015}. The \emph{batch layer} keeps all raw data in a master
dataset that never changes and recomputes views from it on a schedule. The \emph{speed layer} processes
new events as they arrive and covers the time until the next batch run. The \emph{serving layer} merges
both kinds of views for users.

\textbf{Kappa} has only a stream layer \cite{kreps2014}. All data, old and new, is a stream in a log such
as Kafka. There is one processing job. To recompute, a second copy of the job reads the log again from an
old offset, and users switch to it when it catches up.

\subsection{Comparing the options with our requirements}
Table~\ref{tab:lambda-kappa} compares the two options against the requirements in
Table~\ref{tab:requirements}.

\begin{table}[!htb]
\centering\small
\caption{Lambda and Kappa compared for our use case. The last column shows which option fits us better.}
\label{tab:lambda-kappa}
\begin{tabularx}{\textwidth}{@{}P{2.6cm} L L P{1.2cm}@{}}
\toprule
\textbf{Criterion} & \textbf{Lambda} & \textbf{Kappa} & \textbf{Better} \\
\midrule
Latency (R1 to R3) & The speed layer updates every 5 to 10 seconds. & The stream job updates every few seconds. & equal \\
Exact money (R4, R5) & The batch layer reads the complete day after late events and the cost file have arrived. & Windows close at a watermark. Later events are dropped or need extra update logic after results are published. & Lambda \\
Daily cost file & Read directly as a batch input. & Must become a stream and wait in state until a full day of trips is ready. & Lambda \\
Recompute a day (R6) & Rerun the Airflow DAG for that day from Parquet files. & Replay the Kafka log from an old offset with a second copy of the job. & Lambda \\
Storage cost & Cheap compressed Parquet files keep all history. & Kafka must keep weeks or months of events. & Lambda \\
Code & Two paths (stream and batch) that can drift apart. & One code path. & Kappa \\
Operations & More parts: Airflow, a batch job and merge views. & Fewer parts. & Kappa \\
Consistency & Live views are temporary. Batch views become the truth for finished days. & One view, but its accuracy depends on the watermark. & Lambda \\
\bottomrule
\end{tabularx}
\end{table}

\subsection{Our decision: Lambda}
We chose the Lambda architecture for five reasons.
\begin{enumerate}
\item \textbf{Two very different sources.} One source is a continuous stream and the other is one file
per day. The daily file is a natural batch input, and the daily report is a natural batch output.
\item \textbf{Money must be exact.} Vehicles in our simulation sometimes lose signal for 20 to 60 simulated
minutes and then send all the old events at once. In our run, about \LateEventsPerDay{} events per day
arrived more than 10 simulated minutes late. To keep its state small, the speed layer closes each window
soon after it ends, so events that arrive later are lost. The batch layer starts only after the day and a
grace period have ended, so it counts every event, removes repeats over the whole day and joins the complete
cost file.
\item \textbf{Recomputing is cheap and simple.} One simulated day has about \EventsPerDay{} events, which is
only a few megabytes of Parquet. If a partner sends a corrected cost file, we rerun the DAG for that day.
Nothing in Kafka has to be replayed and Kafka can keep a short retention.
\item \textbf{Latency is still low.} The speed layer writes to PostgreSQL every 5 to 10 seconds, so the
live dashboard meets requirements R1 to R3.
\item \textbf{The tools fit.} Airflow is built for scheduled batch jobs, and Spark can run both the
streaming job and the batch job with the same API.
\end{enumerate}

\subsection{The rejected alternative: Kappa}
In a Kappa design, the batch simulator would send its cost rows into a Kafka topic. One Spark streaming job
would keep a full day of trip totals per vehicle in its state and join them with the cost rows when they
arrive. To fix a day, we would replay the log from an old offset. We rejected this design because the join
needs a whole day of state for every vehicle, because late events after the watermark would change money
numbers that were already published, and because replay needs Kafka to keep all history. Kappa does have a
real advantage: one code path and fewer moving parts. If costs arrived as live events, or if the business
only needed live numbers, Kappa would be the better choice.

\subsection{How we reduce the weak points of Lambda}
The main problem of Lambda is keeping two code paths correct and in step. We reduce it in four ways.
\begin{itemize}
\item \textbf{One engine.} Both layers use Spark 4.1 and the same DataFrame API.
\item \textbf{Shared code.} Parsing, validation, the zone lookup, the metric definitions and the profit rules
live in one module (\code{fleet\_transforms.py}) that both layers import. A unit test checks that the Spark
version and the Python version of the profit rules give the same answer.
\item \textbf{A clear hand-over.} For a finished day the serving views show batch results only. The DAG then
deletes speed-layer rows for older days (Figure~\ref{fig:handover}).
\item \textbf{Reconciliation.} After each batch run we compare the speed and batch totals of the day and store
the difference (Section~\ref{sec:reconciliation}).
\end{itemize}

\begin{figure}[!htb]
\centering
% Timeline of the Lambda hand-over from the speed layer to the batch layer
\begin{tikzpicture}[font=\sffamily\scriptsize, x=0.92cm, y=1cm]
% axis
\draw[-{Stealth[length=2mm]}, line width=0.6pt] (0,0) -- (15.8,0) node[right] {time};
\foreach \x in {0, 5, 10, 15} { \draw[line width=0.6pt] (\x,0.08) -- (\x,-0.08); }
\node[below] at (0,-0.08) {00:00};
\node[below] at (5,-0.08) {00:00};
\node[below] at (10,-0.08) {00:00};
\node[below] at (15,-0.08) {00:00};
\node[font=\sffamily\scriptsize\bfseries] at (2.5,-0.45) {day D-2};
\node[font=\sffamily\scriptsize\bfseries] at (7.5,-0.45) {day D-1};
\node[font=\sffamily\scriptsize\bfseries] at (12.5,-0.45) {day D (today)};
% served views
\node[anchor=east, align=right] at (-0.15,0.75) {served\\view};
\fill[green!25, draw=green!50!black, line width=0.4pt] (0.05,0.45) rectangle (4.95,1.05);
\node at (2.5,0.75) {batch view (exact)};
\fill[green!25, draw=green!50!black, line width=0.4pt] (5.05,0.45) rectangle (9.95,1.05);
\node at (7.5,0.75) {batch view (exact)};
\fill[orange!25, draw=orange!60!black, line width=0.4pt] (10.05,0.45) rectangle (13.4,1.05);
\node at (11.72,0.75) {speed view (live)};
% now marker
\draw[dashed, line width=0.6pt] (13.4,-0.05) -- (13.4,1.55);
\node[above] at (13.4,1.55) {now};
% batch run for D-1
\draw[-{Stealth[length=2mm]}, line width=0.6pt] (10.35,1.75) -- (10.35,1.10);
\node[above, align=center] at (10.35,1.75) {cost file for D-1 arrives: the batch run\\replaces the speed rows of D-1};
\end{tikzpicture}

\caption{How the serving layer uses the two layers over time. Finished days come from the batch layer.
Today comes from the speed layer until the batch run for today replaces it after midnight.}
\label{fig:handover}
\end{figure}
